\section*{第16章：子类型化的元理论}
\addcontentsline{toc}{section}{第16章：子类型化的元理论}

\begin{taplsolutionentrybox}
\phantomsection
\label{sol:translator:16.2.1}
\textbf{练习16.2.1}

证明分为两种情形。

\textbf{情形 \textsc{T-Rcd}。}考虑一棵以 \textsc{T-Rcd} 结束的推导。
它的第 \(i\) 个前提为 \(\Gamma\vdash t_i:T_i\)，表示记录中标签 \(l_i\)
对应的字段项 \(t_i\) 具有类型 \(T_i\)。对每个字段分别定义 \(S_i\)：若该
前提的最后一条规则是 \textsc{T-Sub}，就取包容前的类型 \(S_i\)，使
\(\Gamma\vdash t_i:S_i\) 且 \(S_i\subtype T_i\)；若最后一条规则不是
\textsc{T-Sub}，则令 \(S_i=T_i\)，并由自反性得到
\(S_i\subtype T_i\)。先对所有字段应用
\textsc{T-Rcd}，得到
\[
  \Gamma\vdash\{l_i=t_i\}^{i\in1..n}:
  \{l_i:S_i\}^{i\in1..n}
\]
再由记录子类型规则得到
\(\{l_i:S_i\}^{i\in1..n}\subtype
  \{l_i:T_i\}^{i\in1..n}\)，最后在 \textsc{T-Rcd} 之后统一使用一次
\textsc{T-Sub}。

\textbf{情形 \textsc{T-Proj}。}考虑一棵以 \textsc{T-Proj} 结束、且其
前提的最后一条规则是 \textsc{T-Sub} 的推导。设包容前已经得到
\(\Gamma\vdash t:S\)，再由
\(S\subtype\{l_i:T_i\}^{i\in1..n}\) 推出
\(\Gamma\vdash t:\{l_i:T_i\}^{i\in1..n}\)，使 \textsc{T-Proj} 能够
导出 \(\Gamma\vdash t.l_j:T_j\)。由子类型反演，\(S\) 本身也必须是
记录类型，其中包含字段 \(l_j:S_j\)，且 \(S_j\subtype T_j\)。因此，
可以直接从 \(\Gamma\vdash t:S\) 使用 \textsc{T-Proj} 得到
\(\Gamma\vdash t.l_j:S_j\)，再把 \textsc{T-Sub} 作为整个投影推导的
最后一条规则，得到
\(\Gamma\vdash t.l_j:T_j\)。

两种情形都把原先位于某个前提末尾的 \textsc{T-Sub} 越过
\textsc{T-Rcd} 或 \textsc{T-Proj}，使它成为导出相同结论的最后一条规则。

\taplsolutionbacklink{ex:16.2.1}{返回练习16.2.1}
\end{taplsolutionentrybox}

\begin{taplsolutionentrybox}
\phantomsection
\label{sol:translator:16.2.3}
\textbf{练习16.2.3}

取
\[
  s=(\lambda r:\{x:\tapltype{Nat}\}.\,r)\,
      \{x=0,y=1\}
  \qquad\text{和}\qquad
  t=\{x=0,y=1\}
\]
有 \(s\trans t\)。算法类型规则给 \(s\) 的类型是
\(S=\{x:\tapltype{Nat}\}\)。具体地说，抽象中的参数标注使函数具有类型
\[
  \{x:\tapltype{Nat}\}\to\{x:\tapltype{Nat}\}
\]
实参的算法类型虽然是
\(\{x:\tapltype{Nat},y:\tapltype{Nat}\}\)，却是参数类型
\(\{x:\tapltype{Nat}\}\) 的子类型，所以 \textsc{TA-App} 可以使用，并以函数的
结果类型 \(S\) 作为整个应用的类型。

求值时，函数体只是变量 \(r\)，所以一次 \(\beta\) 归约会原样返回实参；参数上的
类型标注只限制函数体可以怎样使用 \(r\)，并不会在运行时删除字段 \(y\)。因此
求值结果
\[
  t=\{x=0,y=1\}
\]
的算法类型为
\[
  T=\{x:\tapltype{Nat},y:\tapltype{Nat}\}
\]
由记录宽度子类型规则，\(T\subtype S\)；反向的 \(S\subtype T\) 不成立。
因此，单步求值使算法算出的最小类型从 \(S\) 变成了更精确的 \(T\)，算法类型
关系不满足“求值前后保持同一类型”这一通常形式的保持性定理。

声明式关系仍然满足通常的保持性。它先给 \(t\) 赋予类型 \(T\)，再由
\(T\subtype S\) 使用 \textsc{T-Sub}，仍可得到 \(\vdash t:S\)。因此这里失效的
只是算法类型的精确相等，类型安全并未受到影响。

\taplsolutionbacklink{ex:16.2.3}{返回练习16.2.3}
\end{taplsolutionentrybox}
